\subsubsection{Transit or Bounce Time}

\paragraph{Normalized Expression}

The transit, or bounce time, is defined as the time for a passing particle to
complete a full orbit in the poloidal plane, and for a trapped particle to
complete half a bounce period. Note that this is possible only in the
approximation of zero banana width. Otherwise, the bounce motion would be no
longer symmetric in the forth and back motions, and both would need to be
acounted for. We define then
\begin{align}
\tau_{b}\left(  \psi\right)   & =\int_{s_{\min}}^{s_{\max}}\frac{ds}{\left|
v_{s}\right|  }\\
& =\int_{s_{\min}}^{s_{\max}}\frac{ds}{\left|  v_{\Vert}\right|  }\frac{B}{
B_{P}}%
\end{align}
where $v_{s}$ is the guiding center velocity along the poloidal field lines,
and $v_{\Vert}$ is its velocity parallel to the magnetic field. $B$ is the
magnitude of the magnetic field, while $B_{P}$ is the magnitude of its
poloidal component. The limits $s_{\min}$ and $s_{\max}$ are defined in
(\ref{bmax}) for passing electrons, and are the positions, along the field
lines, of turning points for trapped electrons.

The differential arc length $ds$ along the poloidal field line is generally
expressed in curvilinear coordinates $\left(  u^{1},u^{2},u^{3}\right)  $ as
(\ref{dl})
\begin{equation}
dl=\sqrt{g_{ij}du^{i}du^{j}}\label{dl2}%
\end{equation}
where the $g_{ij}$ are the metric coefficients, defined in (\ref{met}). In the
$\left(  \psi,\theta,\phi\right)  $ coordinates, the variations $d\psi$ and
$d\phi$ are essentially zero along the poloidal field line. As a consequence,
(\ref{dl2}) becomes
\begin{equation}
dl=\sqrt{g_{22}}d\theta
\end{equation}

The velocity and momentum are related through the relativistic factor
$\gamma\left(  p\right)  $.
\begin{equation}
v=\frac{p}{\gamma m}%
\end{equation}
with
\begin{equation}
\gamma\left(  p\right)  =\sqrt{1+\frac{p^{2}}{m^{2}c^{2}}}%
\end{equation}

Therefore, we have
\begin{equation}
\frac{v_{\Vert}}{v}=\frac{p_{\Vert}}{p}=\xi
\end{equation}
where the pitch-angle cosine $\xi$ was defined in (\ref{pxidef})

We get
\begin{equation}
\tau_{b}\left(  \psi\right)  =\frac{2\pi}{v\left|  \xi_{0}\right|  }%
\int_{\theta_{\min}}^{\theta_{\max}}\frac{d\theta}{2\pi}\sqrt{g_{22}}\frac
{\xi_{0}}{\xi}\frac{B}{B_{P}}%
\end{equation}
where $\xi_{0}$ is the pitch angle cosine at the position $\theta_{0}$ of
minimum $B$-field
\begin{equation}
\theta_{0}\equiv\theta\left(  B=B_{0}\left(  \psi\right)  \right)
\end{equation}
and the limits$\ \theta_{\min}$ and $\theta_{\max}$ will be calculated in the
next subsection.

The bounce time can be normalized as such:
\begin{equation}
\tau_{b}\left(  \psi,\xi_{0}\right)  =\frac{2\pi R_{p}\widetilde{q}\left(
\psi\right)  }{v\left|  \xi_{0}\right|  }\lambda\left(  \psi,\xi_{0}\right)
\end{equation}
with
\begin{equation}
\lambda\left(  \psi,\xi_{0}\right)  =\frac{1}{\widetilde{q}\left(
\psi\right)  }\int_{\theta_{\min}}^{\theta_{\max}}\frac{d\theta}{2\pi}%
\frac{\sqrt{g_{22}}}{R_{p}}\frac{\xi_{0}}{\xi}\frac{B}{B_{P}}\label{lambda0}%
\end{equation}
and
\begin{equation}
\widetilde{q}\left(  \psi\right)  \equiv\int_{0}^{2\pi}\frac{d\theta}{2\pi
}\frac{\sqrt{g_{22}}}{R_{p}}\frac{B}{B_{P}}%
\end{equation}

The bounce time is normalized to the transit time of particles with parallel
momentum only, such that $\lambda\left(  \psi,\pm1\right)  =1$.

The covariant metric element $g_{22}$ is given by (\ref{rel1})-(\ref{met}),
which is in the $\left(  \psi,\theta,\phi\right)  $ system becomes
(\ref{gijpsitheta})
\[
g_{22}=\left|  J\nabla\psi\times\nabla\phi\right|  ^{2}=\frac{r}{\left|
\widehat{\psi}\cdot\widehat{r}\right|  }
\]

Consequently, the normalized bounce time takes the form
\begin{equation}
\lambda\left(  \psi,\xi_{0}\right)  =\frac{1}{\widetilde{q}\left(
\psi\right)  }\int_{\theta_{\min}}^{\theta_{\max}}\frac{d\theta}{2\pi}\frac
{1}{\left|  \widehat{\psi}\cdot\widehat{r}\right|  }\frac{r}{R_{p}}\frac
{B}{B_{P}}\frac{\xi_{0}}{\xi}\label{lambda}%
\end{equation}
with
\begin{equation}
\widetilde{q}\left(  \psi\right)  =\int_{0}^{2\pi}\frac{d\theta}{2\pi}\frac
{1}{\left|  \widehat{\psi}\cdot\widehat{r}\right|  }\frac{r}{R_{p}}\frac
{B}{B_{P}}\label{qtilde}%
\end{equation}

\paragraph{Particle Motion in the Magnetic Field}

The particle motion along the magnetic field lines exhibits one constant of
the motion, the energy (or the total momentum), and an adiabatic invariant,
the magnetic moment. They are given by the equations
\begin{align}
p^{2}  & =p_{\bot}^{2}+p_{\Vert}^{2}\\
\mu & =\frac{p_{\bot}^{2}}{2mB}\nonumber
\end{align}
such that, as a function of the moment component $\left(  p_{\Vert0},p_{\bot
0}\right)  $ at the location $\theta_{0}$ of minimum $B$-field, we have
\begin{align}
p_{\bot}^{2}+p_{\Vert}^{2}  & =p_{\bot0}^{2}+p_{\Vert0}^{2}\label{mom}\\
\frac{p_{\bot}^{2}}{B\left(  \psi,\theta\right)  }  & =\frac{p_{\bot0}^{2}%
}{B_{0}\left(  \psi\right)  }\label{mhu}%
\end{align}

Using the transformation (\ref{pxi}-\ref{pparpperp}) from $\left(  p_{\Vert
},p_{\bot}\right)  $ to $\left(  p,\xi\right)  $, the system (\ref{mom}%
-\ref{mhu}) becomes
\begin{align}
p^{2}  & =p_{0}^{2}\\
\frac{1-\xi^{2}}{B\left(  \psi,\theta\right)  }  & =\frac{1-\xi_{0}^{2}}%
{B_{0}\left(  \psi\right)  }\nonumber
\end{align}

We get an expression for $\xi$ as a function of $\xi_{0}$:
\begin{equation}
\xi\left(  \psi,\theta,\xi_{0}\right)  =\text{sign}\left(  \xi_{0}\right)
\sqrt{1-\Psi\left(  \psi,\theta\right)  \left(  1-\xi_{0}^{2}\right)
}\label{xi}%
\end{equation}
where we define the ratio
\begin{equation}
\Psi\left(  \psi,\theta\right)  \equiv\frac{B\left(  \psi,\theta\right)
}{B_{0}\left(  \psi\right)  }\label{Psidef}%
\end{equation}

The trapping condition is given by $\left|  \xi_{0}\right|  <\xi_{0T}\left(
\psi\right)  $ , where $\xi_{0T}\left(  \psi\right)  $ is the pitch angle,
defined at the minimum $B_{0}\left(  \psi\right)  $ on a given flux-surface,
such that the parallel velocity of the particle vanishes at the maximum
$B_{\max}\left(  \psi\right)  $. An expression for $\xi_{0T}\left(
\psi\right)  $ can then be obtained from (\ref{xi}): setting $\xi\left(
\xi_{0T},B=B_{\max}\left(  \psi\right)  \right)  =0$, we get
\begin{equation}
\xi_{0T}^{2}\left(  \psi\right)  =1-\frac{B_{0}\left(  \psi\right)  }{B_{\max
}\left(  \psi\right)  }\label{xi0t2}%
\end{equation}

The turning points are
\begin{equation}
\theta_{\min}\left(  \psi,\xi_{0}\right)  =\left|
\begin{array}
[c]{cc}%
-\pi & \text{for passing particles}\\
\theta_{T\min} & \text{for trapped particles}%
\end{array}
\right.
\end{equation}%
\begin{equation}
\theta_{\max}\left(  \psi,\xi_{0}\right)  =\left|
\begin{array}
[c]{cc}%
\pi & \text{for passing particles}\\
\theta_{T\max} & \text{for trapped particles}%
\end{array}
\right.
\end{equation}

We can determine the turning angles $\theta_{T\min}\left(  \psi,\xi
_{0}\right)  $ and $\theta_{T\max}\left(  \psi,\xi_{0}\right)  $ as the
position where $\xi\left(  \psi,\theta,\xi_{0}\right)  =0$. At this position,
we have $B=B_{b}\left(  \psi,\xi_{0}\right)  $, where $B_{b}\left(  \psi
,\xi_{0}\right)  $ is then given by (\ref{xi})
\begin{equation}
B_{b}\left(  \psi,\xi_{0}\right)  =\frac{B_{0}\left(  \psi\right)  }{1-\xi
_{0}^{2}}%
\end{equation}
so that
\begin{align}
\theta_{T\min}\left(  \psi,\xi_{0}\right)    & =\theta\left(  B=B_{b}\left|
\theta<\theta_{0}\right.  \right)  \;[2\pi]\label{thetamin}\\
\theta_{T\max}\left(  \psi,\xi_{0}\right)    & =\theta\left(  B=B_{b}\left|
\theta>\theta_{0}\right.  \right)  \;[2\pi]\nonumber
\end{align}

\paragraph{Calculation of $\lambda\left(  \psi,\xi_{0}\right)  $ From the
Output Data of Equilibrium Codes}

The numerical calculation of $\lambda\left(  \psi,\xi_{0}\right)  $ can be
carried from the output of equilibrium codes such as HELENA. These data are
assumed to be the parametrization of the flux-surfaces $R\left(  \psi
,\theta\right)  $ and $Z\left(  \psi,\theta\right)  $, and the three
components of the magnetic field $B_{R}\left(  \psi,\theta\right)  $,
$B_{Z}\left(  \psi,\theta\right)  $ and $B_{\phi}\left(  \psi,\theta\right)
$. From these components we derive directly the toroidal and poloidal
components of the field, as well as the total field:
\begin{align}
B_{T}\left(  \psi,\theta\right)    & =\left|  B_{\phi}\left(  \psi
,\theta\right)  \right|  \nonumber\\
B_{P}\left(  \psi,\theta\right)    & =\sqrt{B_{R}^{2}\left(  \psi
,\theta\right)  +B_{Z}^{2}\left(  \psi,\theta\right)  }\nonumber\\
B\left(  \psi,\theta\right)    & =\sqrt{B_{T}^{2}\left(  \psi,\theta\right)
+B_{P}^{2}\left(  \psi,\theta\right)  }\label{Btot}%
\end{align}
and also%
\begin{align}
R_{p}  & =R\left(  0,\theta\right)  \\
Z_{p}  & =Z\left(  0,\theta\right)
\end{align}

We also have an expression for $r$
\begin{equation}
r\left(  \psi,\theta\right)  =\sqrt{\left(  R\left(  \psi,\theta\right)
-R_{p}\right)  ^{2}+\left(  Z\left(  \psi,\theta\right)  -Z_{p}\right)  ^{2}%
}\label{r}%
\end{equation}
and, using
\begin{align}
\widehat{r}  & =\left(  \widehat{r}\cdot\widehat{R}\right)  \widehat
{R}+\left(  \widehat{r}\cdot\widehat{Z}\right)  \widehat{Z}\nonumber\\
& =\left(  \frac{R\left(  \psi,\theta\right)  -R_{p}}{r}\right)  \widehat
{R}+\left(  \frac{Z\left(  \psi,\theta\right)  -Z_{p}}{r}\right)  \widehat{Z}%
\end{align}
we get an expression for the scalar product
\begin{equation}
\widehat{\psi}\cdot\widehat{r}=\frac{\left(  \nabla\psi\cdot\widehat
{R}\right)  \left(  R\left(  \psi,\theta\right)  -R_{p}\right)  +\left(
\nabla\psi\cdot\widehat{Z}\right)  \left(  Z\left(  \psi,\theta\right)
-Z_{p}\right)  }{r\left|  \nabla\psi\right|  }\label{psidotr}%
\end{equation}

In a toroidal axisymmetric geometry, the magnetic field can be expressed
generally as
\begin{equation}
\mathbf{B}=I\left(  \psi\right)  \nabla\phi+\nabla\psi\times\nabla
\phi\label{B}%
\end{equation}
so that
\begin{align}
B_{T}  & =\left|  I\left(  \psi\right)  \right|  \left|  \nabla\phi\right|
=\frac{\left|  I\left(  \psi\right)  \right|  }{R}\\
B_{P}  & =\left|  \nabla\psi\right|  \left|  \nabla\phi\right|  =\frac{\left|
\nabla\psi\right|  }{R}\nonumber
\end{align}

We also have
\begin{align}
\mathbf{B}_{T}  & =I\left(  \psi\right)  \nabla\phi=B_{\phi}\widehat{\phi
}\label{Bp}\\
\mathbf{B}_{P}  & =\nabla\psi\times\nabla\phi=-B_{P}\widehat{s}\nonumber
\end{align}
and therefore
\begin{align}
\nabla\phi\times\mathbf{B}_{P}  & =\nabla\phi\times\left(  \nabla\psi
\times\nabla\phi\right) \nonumber\\
& =\left|  \nabla\phi\right|  ^{2}\nabla\psi
\end{align}
so that
\begin{equation}
\nabla\psi=\frac{\nabla\phi\times\mathbf{B}_{P}}{\left|  \nabla\phi\right|
^{2}}=R\widehat{\phi}\times\mathbf{B}_{P}%
\end{equation}
and we have the projections
\begin{align}
\left(  \nabla\psi\cdot\widehat{R}\right)   & =R\widehat{R}\cdot\widehat{
\phi}\times\mathbf{B}_{P}=-RB_{Z}\\
\left(  \nabla\psi\cdot\widehat{Z}\right)   & =R\widehat{Z}\cdot\widehat{
\phi}\times\mathbf{B}_{P}=RB_{R}\nonumber
\end{align}

Finally, the expression for the normalized bounce time becomes
\begin{equation}
\lambda\left(  \psi,\xi_{0}\right)  =\frac{1}{\widetilde{q}\left(
\psi\right)  }\int_{\theta_{\min}}^{\theta_{\max}}\frac{d\theta}{2\pi}%
\frac{B\left[  \left(  R-R_{p}\right)  ^{2}+\left(  Z-Z_{p}\right)
^{2}\right]  }{R_{p}\left|  B_{R}\left(  Z-Z_{p}\right)  -B_{Z}\left(
R-R_{p}\right)  \right|  }\frac{\xi_{0}}{\xi}\label{lambda1}%
\end{equation}
with
\begin{equation}
\widetilde{q}\left(  \psi\right)  =\int_{0}^{2\pi}\frac{d\theta}{2\pi}%
\frac{B\left[  \left(  R-R_{p}\right)  ^{2}+\left(  Z-Z_{p}\right)
^{2}\right]  }{R_{p}\left|  B_{R}\left(  Z-Z_{p}\right)  -B_{Z}\left(
R-R_{p}\right)  \right|  }%
\end{equation}
where $R,Z,B_{R},B_{Z}$ and $B$ are functions of $\left(  \psi,\theta\right)
$, and $\xi$ is a function of $\left(  \psi,\theta,\xi_{0}\right)  $ given by
(\ref{xi})
\begin{equation}
\xi\left(  \psi,\theta,\xi_{0}\right)  =\text{sign}\left(  \xi_{0}\right)
\sqrt{1-\frac{B\left(  \psi,\theta\right)  }{B_{0}\left(  \psi\right)
}\left(  1-\xi_{0}^{2}\right)  }%
\end{equation}

This will be used to calculate $\lambda$ numerically.

\paragraph{Safety Factor $q\left(  \psi\right)  $}

The (averaged) safety factor is generally defined by [callen 11.5.6]
\begin{equation}
q\left(  \psi\right)  =\frac{I\left(  \psi\right)  }{4\pi^{2}}\frac{\delta
V}{\delta\psi}\left\langle R^{-2}\right\rangle
\end{equation}
where $V$ is the volume enclosed by a flux-surface and $\left\langle
\;\right\rangle $ denotes the flux-surface average.

It can be expressed as
\begin{equation}
q\left(  \psi\right)  =I\left(  \psi\right)  \int_{0}^{2\pi}\frac{d\theta}{
2\pi}\int_{0}^{2\pi}\frac{d\phi}{2\pi}\frac{J}{R^{2}}%
\end{equation}
where the Jacobian $J$ is given by (\ref{jacpsitheta})
\begin{align}
J  & =\left|  \nabla\psi\times\nabla\theta\cdot\nabla\phi\right|
^{-1}\nonumber\\
& =\frac{Rr}{\left|  \nabla\psi\right|  \left|  \widehat{\psi}\cdot\widehat{
r}\right|  }\nonumber\\
& =\frac{r}{B_{P}}\frac{1}{\left|  \widehat{\psi}\cdot\widehat{r}\right|  }%
\end{align}
where we used (\ref{Bp}).

We obtain
\begin{equation}
q\left(  \psi\right)  =I\left(  \psi\right)  \int_{0}^{2\pi}\frac{d\theta}{
2\pi}\frac{1}{\left|  \widehat{\psi}\cdot\widehat{r}\right|  }\frac{r}{
B_{P}R^{2}}%
\end{equation}
and, using (\ref{B}), we finally have
\begin{equation}
q\left(  \psi\right)  =\int_{0}^{2\pi}\frac{d\theta}{2\pi}\frac{1}{\left|
\widehat{\psi}\cdot\widehat{r}\right|  }\frac{r}{R}\frac{B_{T}}{B_{P}%
}\label{qfactor}%
\end{equation}

The expression of $q\left(  \psi\right)  $ and its relation to $\widetilde
{q}\left(  \psi\right)  $ in the simplified case of circular concentric
flux-surfaces will be addressed in Subsection \ref{circonc}.

Using (\ref{r}) and (\ref{psidotr}), we find the expression
\begin{equation}
q\left(  \psi\right)  =\int_{0}^{2\pi}\frac{d\theta}{2\pi}\frac{\left[
\left(  R-R_{p}\right)  ^{2}+\left(  Z-Z_{p}\right)  ^{2}\right]  B_{T}%
}{R\left|  B_{R}\left(  Z-Z_{p}\right)  -B_{Z}\left(  R-R_{p}\right)  \right|
}\label{qfactor2}%
\end{equation}

\subsubsection{Toroidal Extent of Banana Orbits}

We are interested in calculating the toroidal extent of banana orbits, that
is, the toroidal angle corresponding to the path done by a trapped particle
between two turning points. It is given by
\begin{equation}
\Delta\phi=\phi_{\max}-\phi_{\min}=\int_{\phi_{\min}}^{\phi_{\max}}d\phi
\end{equation}
and can be expressed as a function of the length element along the path, using
(\ref{dlphi})
\begin{equation}
\Delta\phi=\int_{l\left(  \phi_{\min}\right)  }^{l\left(  \phi_{\max}\right)
}dl\left(  \phi\right)  \frac{d\phi}{dl\left(  \phi\right)  }=\int_{l\left(
\phi_{\min}\right)  }^{l\left(  \phi_{\max}\right)  }\frac{dl\left(
\phi\right)  }{R}%
\end{equation}

The poloidal and toroidal elements are related through the local angle of the
magnetic field,
\begin{equation}
\frac{dl\left(  \phi\right)  }{dl\left(  \theta\right)  }=\frac{B_{T}}{B_{P}}%
\end{equation}
so that
\begin{equation}
\Delta\phi=\int_{l\left(  \phi_{\min}\right)  }^{l\left(  \phi_{\max}\right)
}dl\left(  \phi\right)  \frac{d\phi}{dl\left(  \phi\right)  }=\int_{l\left(
\theta_{\min}\right)  }^{l\left(  \theta_{\max}\right)  }\frac{1}{R}%
\frac{B_{T}}{B_{P}}dl\left(  \theta\right)
\end{equation}

Using (\ref{dltheta}), we get
\begin{equation}
\Delta\phi=\int_{\theta_{\min}}^{\theta_{\max}}d\theta\;\frac{1}{\left|
\widehat{\psi}\cdot\widehat{r}\right|  }\frac{r}{R}\frac{B_{T}}{B_{P}}%
\end{equation}

Defining the integral
\begin{equation}
q_{T}\left(  \psi,\xi_{0}\right)  =\int_{\theta_{\min}}^{\theta_{\max}}%
\frac{d\theta}{2\pi}\;\frac{1}{\left|  \widehat{\psi}\cdot\widehat{r} \right|
}\frac{r}{R}\frac{B_{T}}{B_{P}}\label{qtdef}%
\end{equation}
we find that the toroidal extent of banana orbits is
\begin{equation}
\frac{\Delta\phi}{2\pi}=q_{T}\left(  \psi,\xi_{0}\right)
\end{equation}

Note that at the trapped/passing limit, we have
\begin{equation}
\lim_{\xi_{0}\rightarrow\xi_{0T}}\frac{\Delta\phi}{2\pi}=q_{T}\left(  \psi
,\xi_{0T}\right)  =q\left(  \psi\right)
\end{equation}
Therefore, we retrieve the interpretation of the safety factors, which is the
number of toroidal rotations $\Delta\phi/2\pi$ for one poloidal rotation.

\subsubsection{Bounce Average}

In order to reduce the dimension of kinetic equations, it is important to
define an average over the poloidal motion, which anihilates the term that
accounts for the time evolution of the variations of the distribution function
along the field lines. The natural average is
\begin{equation}
\left\{  \mathcal{A}\right\}  =\frac{1}{\tau_{b}}\left[  \frac{1}{2}%
\sum_{\sigma}\right]  _{T}\int_{s_{\min}}^{s_{\max}}\frac{ds}{\left|
v_{s}\right|  }\mathcal{A}\label{bav0}%
\end{equation}
where $\sigma=$sign$\left(  v_{\Vert}\right)  $ and the sum over $\sigma$
applies to trapped particles only.

It can be rewritten in terms of the normalized bounce time $\lambda$
(\ref{lambda0})
\begin{equation}
\left\{  \mathcal{A}\right\}  =\frac{1}{\lambda\widetilde{q}}\left[  \frac
{1}{2}\sum_{\sigma}\right]  _{T}\int_{\theta_{\min}}^{\theta_{\max}}%
\frac{d\theta}{2\pi}\frac{\sqrt{g_{22}}}{R_{p}}\frac{B}{B_{P}}\frac{\xi_{0}%
}{\xi}\mathcal{A}\label{bav}%
\end{equation}
or, expressing the metric coefficient,
\begin{equation}
\left\{  \mathcal{A}\right\}  =\frac{1}{\lambda\widetilde{q}}\left[  \frac
{1}{2}\sum_{\sigma}\right]  _{T}\int_{\theta_{\min}}^{\theta_{\max}}%
\frac{d\theta}{2\pi}\frac{1}{\left|  \widehat{\psi}\cdot\widehat{r}\right|
}\frac{r}{R_{p}}\frac{B}{B_{P}}\frac{\xi_{0}}{\xi}\mathcal{A}\label{bav2}%
\end{equation}

Another expression uses the output data from equilibrium codes. Following the
work in the previous section, we find
\begin{equation}
\left\{  \mathcal{A}\right\}  =\frac{1}{\lambda\widetilde{q}}\left[  \frac
{1}{2}\sum_{\sigma}\right]  _{T}\int_{\theta_{\min}}^{\theta_{\max}}%
\frac{d\theta}{2\pi}\frac{B\left[  \left(  R-R_{p}\right)  ^{2}+\left(
Z-Z_{p}\right)  ^{2}\right]  }{R_{p}\left|  B_{R}\left(  Z-Z_{p}\right)
-B_{Z}\left(  R-R_{p}\right)  \right|  }\frac{\xi_{0}}{\xi}\mathcal{A}%
\end{equation}
or explicitely
\begin{align}
\left\{  \mathcal{A}\right\}    & =\left[  \int_{\theta_{\min}}^{\theta_{\max
}}\frac{d\theta}{2\pi}\frac{B\left[  \left(  R-R_{p}\right)  ^{2}+\left(
Z-Z_{p}\right)  ^{2}\right]  }{R_{p}\left|  B_{R}\left(  Z-Z_{p}\right)
-B_{Z}\left(  R-R_{p}\right)  \right|  }\frac{\xi_{0}}{\xi}\right]
^{-1}\times\\
& \;\left[  \frac{1}{2}\sum_{\sigma}\right]  _{T}\int_{\theta_{\min}}%
^{\theta_{\max}}\frac{d\theta}{2\pi}\frac{B\left[  \left(  R-R_{p}\right)
^{2}+\left(  Z-Z_{p}\right)  ^{2}\right]  }{R_{p}\left|  B_{R}\left(
Z-Z_{p}\right)  -B_{Z}\left(  R-R_{p}\right)  \right|  }\frac{\xi_{0}}{\xi
}\mathcal{A}%
\end{align}

The bounce averaging of momentum-space operators in the kinetic equations
leads to coefficients with similar structure, denoted $\lambda_{k,l,m}$ and
$\overline{\lambda}_{k,l,m}$, which we define as
\begin{equation}
\left\{  \left(  \frac{\xi\left(  \psi,\theta,\xi_{0}\right)  }{\xi_{0}%
}\right)  ^{k}\Psi^{l}\left(  \psi,\theta\right)  \left(  \frac{R_{0}\left(
\psi\right)  }{R\left(  \psi,\theta\right)  }\right)  ^{m}\right\}
=\frac{\lambda_{k,l,m}\left(  \psi,\xi_{0}\right)  }{\lambda\left(  \psi
,\xi_{0}\right)  }\label{lkm}%
\end{equation}
and
\begin{equation}
\sigma\left\{  \sigma\left(  \frac{\xi\left(  \psi,\theta,\xi_{0}\right)
}{\xi_{0}}\right)  ^{k}\Psi^{l}\left(  \psi,\theta\right)  \left(  \frac
{R_{0}\left(  \psi\right)  }{R\left(  \psi,\theta\right)  }\right)
^{m}\right\}  =\frac{\overline{\lambda}_{k,l,m}\left(  \psi,\xi_{0}\right)
}{\lambda\left(  \psi,\xi_{0}\right)  }\label{lkmb}%
\end{equation}
where%
\begin{equation}
R_{0}\left(  \psi\right)  \equiv R\left(  \psi,\theta_{0}\right)
\end{equation}

Note that $\lambda_{0,0,0}=\lambda$. In addition, we see that
\begin{equation}
\overline{\lambda}_{k,l,m}=\left|
\begin{array}
[c]{ll}%
\lambda_{k,l,m} & \text{ for passing particles}\\
0 & \text{ for trapped particles}%
\end{array}
\right.
\end{equation}
